%============================================================
% RUNNER: BUKU KERJA PRAKTIK
% Compile: make buku
%
% Struktur halaman (prelim = romawi, mengikuti Pasal 4.2.1):
%  i    Sampul (nomor tidak ditampilkan)
%  ii   Halaman kosong
%  iii  Halaman Judul Dalam (tanpa nomor, Pasal 4.2.2 4b)
%  iv   Halaman kosong
%  v    Lembar Pengesahan Internal / PENS (tanpa nomor)
%  vi   Halaman kosong
%  vii  Lembar Pengesahan Instansi / Perusahaan (tanpa nomor)
%  viii Halaman kosong
%  ix   Abstrak (nomor mulai tampil)
%  ...  Daftar Isi, Kata Pengantar, Daftar Gambar, Daftar Tabel,
%       Daftar Lampiran (setiap bagian mulai halaman ganjil)
%  ---  BAB 1: angka Arab mulai 1
%============================================================

\documentclass[10pt, a5paper, twoside, openright, fleqn]{report}

%============================================================
% VARIABEL IDENTITAS DOKUMEN KERJA PRAKTIK
% Edit bagian ini sesuai dengan data kerja praktik Anda.
%============================================================

% -- Gaya Header & Footer (Pilih salah satu) -------------------
% 1 = Simple (Laporan kerja praktik, tengah)
% 2 = PENS (Logo PENS, Nama Perusahaan, dll)
\newcommand{\StyleHeaderFooter}{1}

% -- Identitas Mahasiswa 1 (wajib diisi) ----------------------
\newcommand{\NamaMahasiswaSatu}{Eka Danis Saputri}
\newcommand{\NRPSatu}{3123600036}

% -- Identitas Mahasiswa 2 (kosongkan jika tidak ada) ---------
\newcommand{\NamaMahasiswaDua}{Muhammad Satrio Asadel}
\newcommand{\NRPDua}{3123600056}

% -- Identitas Mahasiswa 3 (kosongkan jika tidak ada) ---------
\newcommand{\NamaMahasiswaTiga}{}
\newcommand{\NRPTiga}{}

% -- Judul Laporan Kerja Praktik ------------------------------
\newcommand{\JudulKP}{Pengembangan Sistem Analitik Prediktif, Evaluasi Keusangan Aset, dan Manajemen Konfigurasi Gardu pada Aplikasi Real-Time Monitoring ``Voltra'' PT. Sinergi Dimensi Informatika}

% -- Tempat Pelaksanaan KP ------------------------------------
\newcommand{\NamaPerusahaan}{PT. Sinergi Dimensi Informatika}
% Nama singkat untuk header Gaya 2 (kosongkan untuk memakai nama lengkap)
\newcommand{\NamaPerusahaanSingkat}{PT. SINDIKA}
\newcommand{\AlamatPerusahaan}{JAPFA Indoland Center (Japfa Tower 2), Lantai 12, Jalan Panglima Sudirman Kav. 66-68, Surabaya, Jawa Timur 60271}

% -- Dosen Pembimbing (sertakan gelar akademis) ---------------
\newcommand{\NamaPembimbing}{Umi Sa'adah, S.Kom, M.Kom}
\newcommand{\NIPPembimbing}{197404162000032003}

% -- Pembimbing Lapangan (dari instansi/perusahaan) -----------
\newcommand{\NamaPembimbingLapangan}{Willy Achmat Fauzi, S.T., M.T.}
\newcommand{\JabatanPembimbingLapangan}{Direktur PT. Sinergi Dimensi Informatika}

% -- Koordinator Kerja Praktik --------------------------------
\newcommand{\NamaKoordinatorKP}{Desy Intan Permatasari, S.Kom., M.Kom.}
\newcommand{\NIPKoordinatorKP}{198603232015042004}

% -- Ketua Program Studi (sertakan gelar akademis) ------------
\newcommand{\NamaKetuaProdi}{Yanuar Risah Prayogi, S.Kom., M.T.}
\newcommand{\NIPKetuaProdi}{198801062019031009}

% -- Program Studi & Institusi --------------------------------
% Jenjang: D3, D4, atau Sarjana Terapan
\newcommand{\JenjangProdi}{D4}
\newcommand{\NamaProdi}{Teknik Informatika}
\newcommand{\NamaDepartemen}{Departemen Teknik Informatika Dan Komputer}

% -- Periode Pelaksanaan KP -----------------------------------
% Contoh: 13 Januari 2025 -- 06 Juni 2025
\newcommand{\PeriodeKP}{5 Januari 2026 -- 22 Mei 2026}

% -- Tahun ----------------------------------------------------
\newcommand{\TahunKP}{2026}

% -- Tanggal Pengesahan ---------------------------------------
\newcommand{\TanggalPengesahan}{Surabaya, 1 Juli 2026}

%============================================================
% TOGGLES: Boolean Switch Opsional Buku KP
% Default: semua FALSE. Aktifkan di main_buku.tex
% (setelah %============================================================
% TOGGLES: Boolean Switch Opsional Buku KP
% Default: semua FALSE. Aktifkan di main_buku.tex
% (setelah %============================================================
% TOGGLES: Boolean Switch Opsional Buku KP
% Default: semua FALSE. Aktifkan di main_buku.tex
% (setelah \input{config/toggles}).
%============================================================

% Sisipkan Lembar Pengesahan Internal (PENS) dalam bentuk PDF
% yang sudah ditandatangani, menggantikan halaman hasil LaTeX.
\newif\ifpengesahaninternalpdf  \pengesahaninternalpdffalse
\newcommand{\FilePengesahanInternalPDF}{pages/pengesahan_internal_ttd.pdf}

% Sisipkan Lembar Pengesahan Instansi (perusahaan) dalam bentuk
% PDF yang sudah ditandatangani & distempel perusahaan.
\newif\ifpengesahaninstansipdf  \pengesahaninstansipdffalse
\newcommand{\FilePengesahanInstansiPDF}{pages/pengesahan_instansi_ttd.pdf}
).
%============================================================

% Sisipkan Lembar Pengesahan Internal (PENS) dalam bentuk PDF
% yang sudah ditandatangani, menggantikan halaman hasil LaTeX.
\newif\ifpengesahaninternalpdf  \pengesahaninternalpdffalse
\newcommand{\FilePengesahanInternalPDF}{pages/pengesahan_internal_ttd.pdf}

% Sisipkan Lembar Pengesahan Instansi (perusahaan) dalam bentuk
% PDF yang sudah ditandatangani & distempel perusahaan.
\newif\ifpengesahaninstansipdf  \pengesahaninstansipdffalse
\newcommand{\FilePengesahanInstansiPDF}{pages/pengesahan_instansi_ttd.pdf}
).
%============================================================

% Sisipkan Lembar Pengesahan Internal (PENS) dalam bentuk PDF
% yang sudah ditandatangani, menggantikan halaman hasil LaTeX.
\newif\ifpengesahaninternalpdf  \pengesahaninternalpdffalse
\newcommand{\FilePengesahanInternalPDF}{pages/pengesahan_internal_ttd.pdf}

% Sisipkan Lembar Pengesahan Instansi (perusahaan) dalam bentuk
% PDF yang sudah ditandatangani & distempel perusahaan.
\newif\ifpengesahaninstansipdf  \pengesahaninstansipdffalse
\newcommand{\FilePengesahanInstansiPDF}{pages/pengesahan_instansi_ttd.pdf}


% -- Opsi (uncomment untuk mengaktifkan) ----------------------
% \pengesahaninternalpdftrue  % Pakai PDF Lembar Pengesahan Internal bertanda tangan
%                             % (pages/pengesahan_internal_ttd.pdf)
% \pengesahaninstansipdftrue  % Pakai PDF Lembar Pengesahan Instansi bertanda tangan & stempel
%                             % (pages/pengesahan_instansi_ttd.pdf)

%============================================================
% PACKAGES
% Semua package yang digunakan dalam dokumen ini.
%============================================================

\usepackage[T1]{fontenc}
\usepackage[utf8]{inputenc}
\usepackage{mathptmx}    % Font Times New Roman (teks + math)
\usepackage{setspace}    % Pengaturan spasi baris
\usepackage{geometry}    % Pengaturan margin halaman
\usepackage{graphicx}    % Penyisipan gambar
\usepackage{caption}     % Kustomisasi caption gambar/tabel
\usepackage{float}       % Kontrol posisi float [H]
\usepackage{tabularx}    % Tabel dengan lebar kolom fleksibel
\usepackage{booktabs}    % Tabel bergaya profesional
\usepackage{array}       % Definisi kolom tabel kustom
\usepackage{amsmath}     % Persamaan matematika
\usepackage{amssymb}     % Simbol matematika tambahan
\usepackage{xcolor}      % Definisi warna
\usepackage{colortbl}    % Warna sel tabel (\rowcolor)
\usepackage{tikz}        % Grafis vektor (watermark, dll.)
\usepackage{fancyhdr}    % Header dan footer kustom
% \usepackage{emptypage}  % DINONAKTIFKAN: konflik dengan kpruler di halaman kosong   % Halaman kosong tanpa header/footer
\usepackage{tocloft}     % Kustomisasi Daftar Isi/Gambar/Tabel
\usepackage{titlesec}    % Format judul bab dan subbab
\usepackage{indentfirst} % Indentasi pada paragraf pertama
\usepackage{parskip}     % Pengaturan jarak antar paragraf
\usepackage[round]{natbib}  % sitasi nama-tahun: \citep -> (Nama, 2021); daftar pustaka urut abjad (apalike)
\usepackage{url}         % URL pada daftar pustaka
\usepackage{hyperref}    % Hyperlink internal/eksternal
\usepackage{ragged2e}    % Perataan teks fleksibel
\usepackage{microtype}   % Perbaikan tipografi mikro
\usepackage{enumitem}    % Kustomisasi list/enumerate
\usepackage{pdfpages}    % Menyisipkan halaman PDF eksternal
\usepackage{etoolbox}    % Utilitas pemrograman LaTeX
\usepackage{tabularray}  % Tabel modern (tblr, longtblr)
\usepackage{eso-pic}     % Untuk kontrol watermark background per halaman
\usepackage[indonesian]{babel} % Aturan hyphenation/pemenggalan kata Bahasa Indonesia

%============================================================
% UKURAN FONT
% Semua ukuran font template ada di sini, satu macro per
% peruntukan. Ubah satu baris, berlaku di semua halaman yang
% memakai peran tersebut. Ukuran yang sama tetapi peruntukannya
% berbeda tetap punya macro sendiri, agar bisa diubah terpisah.
%
% Format: \fontsize{ukuran huruf}{jarak baris}\selectfont (pt)
% Di-input setelah config/packages.tex dan sebelum
% config/formatting.tex di main_buku.tex.
%============================================================

% -- Judul Isi Dokumen ----------------------------------------
\newcommand{\FontJudulBab}{\fontsize{12}{15}\selectfont}           % BAB n + judul bab
\newcommand{\FontJudulSubbab}{\fontsize{10}{13}\selectfont}        % 1.1 (section)
\newcommand{\FontJudulSubSubbab}{\fontsize{10}{13}\selectfont}     % 1.1.1 (subsection)
\newcommand{\FontJudulSubSubSubbab}{\fontsize{10}{13}\selectfont}  % tanpa nomor (subsubsection)
\newcommand{\FontJudulDaftar}{\fontsize{12}{15}\selectfont}        % DAFTAR ISI / GAMBAR / TABEL
\newcommand{\FontJudulHalaman}{\fontsize{12}{15}\selectfont}       % ABSTRAK, KATA PENGANTAR, LAMPIRAN
\newcommand{\FontKutipanPanjang}{\fontsize{9}{11}\selectfont}      % lingkungan kutipanpanjang

% -- Caption & Sumber Gambar/Tabel (BIMA.PENS.KP Rev.02) ------
% Judul gambar 10pt bold, judul tabel 10pt normal: diatur di
% config/formatting.tex lewat \captionsetup. Sumber 8pt:
\newcommand{\FontSumber}{\fontsize{8}{10}\selectfont}              % "Sumber: ..." di bawah gambar/tabel

% -- Header & Footer (pilih gaya di config/variables.tex) -----
\newcommand{\FontHeaderFooterSimple}{\fontsize{10}{12}\selectfont} % Gaya 1: Simple
\newcommand{\FontHeaderFooterPENS}{\fontsize{8}{10}\selectfont}    % Gaya 2: PENS

% -- Cover (formats/cover_buku.tex) ---------------------------
\newcommand{\FontCoverLabel}{\fontsize{10}{13}\selectfont}         % "KERJA PRAKTIK"
\newcommand{\FontCoverJudul}{\fontsize{10}{13}\selectfont}         % judul KP
\newcommand{\FontCoverIsi}{\fontsize{10}{13}\selectfont}           % nama, NRP, pembimbing, prodi

% -- Halaman Judul Dalam (formats/halaman_judul.tex) ----------
\newcommand{\FontJudulDalamLabel}{\fontsize{10}{13}\selectfont}    % "LAPORAN KERJA PRAKTIK"
\newcommand{\FontJudulDalamJudul}{\fontsize{10}{13}\selectfont}    % judul KP
\newcommand{\FontJudulDalamIsi}{\fontsize{10}{13}\selectfont}      % periode, nama, NRP, prodi

% -- Lembar Pengesahan Internal & Instansi --------------------
\newcommand{\FontPengesahanJudul}{\fontsize{10}{13}\selectfont}    % blok judul "LEMBAR PENGESAHAN ..."
\newcommand{\FontPengesahanIsi}{\fontsize{10}{13}\selectfont}      % penyusun, tanda tangan, NIP
         % ukuran font per peruntukan
%============================================================
% FORMATTING: BASE LAYER (SHARED)
% Aturan yang 100% sama di semua dokumen KP.
% Aturan per-dokumen (geometry, spacing) ada di formats/format_buku.tex
% yang di-input SETELAH file ini di main_buku.tex.
%============================================================

% -- Warna Institusi ------------------------------------------
\definecolor{pensorange}{RGB}{255, 149, 0}   % Orange KP (#FF9500)
\definecolor{penskuning}{RGB}{255, 193, 7} % Kuning/emas garis cover

% -- Nonaktifkan Hyphenation (pemenggalan kata) ----------------------
% Mencegah LaTeX memotong kata di akhir baris (mis. "Elek-tronika").
% tolerance=9999 + emergencystretch memberi ruang fleksibel agar
% LaTeX tidak perlu Overfull saat hyphenation dinonaktifkan.
\hyphenpenalty=10000
\exhyphenpenalty=10000
\tolerance=9999
\emergencystretch=3em

% -- Lebar Judul di Cover & Halaman Judul (harus sama) ---------
% 95mm = 0.92 * textwidth_normal (0.92 * 103mm = 94.76mm)
\newlength{\LebarJudulKP}
\setlength{\LebarJudulKP}{90mm}  % Silahkan diganti sesuai keinginan

% -- Indentasi & Paragraf -------------------------------------
% \parindent TIDAK diset di sini -- diset di formats/format_buku.tex
\setlength{\parskip}{0pt}

% -- Persamaan Matematika -------------------------------------
\setlength{\mathindent}{10mm}
\setlength{\abovedisplayskip}{12pt}
\setlength{\belowdisplayskip}{12pt}
\setlength{\abovedisplayshortskip}{12pt}
\setlength{\belowdisplayshortskip}{12pt}

% Nomor persamaan dicetak tebal-miring
\makeatletter
\renewcommand{\tagform@}[1]{\maketag@@@{\bfseries\itshape(#1)\@@italiccorr}}

% -- Macro: render fontsize tetap, cuma shrink kalau kepanjangan --
% Usage: \fitline{<max-width>}{<content>}
\newcommand{\fitline}[2]{%
  \setbox0=\hbox{{#2}}%
  \ifdim\wd0>#1
    \resizebox{#1}{!}{{#2}}%
  \else
    {#2}%
  \fi
}
\makeatother

% -- Format Judul Bab -----------------------------------------
% Nomor bab angka Arab: BAB 1, BAB 2, ... (Pasal 4.2.2 no. 3b)
\renewcommand{\thesection}{\arabic{chapter}.\arabic{section}}

% Mulai dari batas margin atas (spacing before = 0pt)
% Judul bab 1 spasi (12pt) di bawah tulisan BAB
% Kalimat awal mulai 2 spasi (24pt) di bawah judul bab
\titleformat{\chapter}[display]
{\centering\FontJudulBab\bfseries}
{\MakeUppercase{\chaptertitlename} \thechapter}
{0pt}
{\MakeUppercase}
\titlespacing*{\chapter}{0pt}{0pt}{12pt}

% -- Auto-selaraskan indentasi paragraf pertama dengan huruf judul (bukan nomor) --
\newsavebox{\KotakLabelJudul}
\newdimen\LebarLabelJudul
\newcommand{\SelaraskanIndentAsli}[1]{%
  \sbox{\KotakLabelJudul}{#1}%
  \global\LebarLabelJudul=\wd\KotakLabelJudul
  \global\advance\LebarLabelJudul by 1em
  \xdef\ParindentAsliTersimpan{\the\parindent}%
  \global\parindent=\LebarLabelJudul
  \global\everypar{\parindent=\ParindentAsliTersimpan\relax\global\everypar{}}%
}

% -- Format Judul Subbab (Section) ----------------------------
\titleformat{\section}
{\FontJudulSubbab\bfseries}
{\thesection}
{1em}
{}
[\SelaraskanIndentAsli{\thesection}]
\titlespacing*{\section}{0pt}{10pt}{3pt}

% -- Format Judul Sub-subbab (Subsection) ---------------------
\titleformat{\subsection}
{\FontJudulSubSubbab\bfseries}
{\thesubsection}
{1em}
{}
[\SelaraskanIndentAsli{\thesubsection}]
\titlespacing*{\subsection}{0pt}{8pt}{3pt}

% -- Format Judul Sub-sub-subbab (Subsubsection) --------------
% Tanpa nomor (secnumdepth=2) - heading ringan, label dikosongkan.
\titleformat{\subsubsection}
{\FontJudulSubSubSubbab\bfseries}
{}
{0pt}
{}
\titlespacing*{\subsubsection}{\parindent}{3pt}{2pt}

% -- Batas Penomoran dan Daftar Isi ---------------------------
% 0 = Chapter, 1 = Section, 2 = Subsection
\setcounter{secnumdepth}{2}
\setcounter{tocdepth}{2}

% -- Caption Gambar (Pasal 4.2.2 no. 6) -----------------------
% "Gambar 2.1: Judul gambar", label & judul bold (contoh pedoman
% hal. 26), di bawah gambar, menjorok 1 cm dari margin kiri-kanan,
% baris lanjutan sejajar judul. Sumber 8pt ditulis setelah caption.
\captionsetup[figure]{
  labelfont       = {bf},
  textfont        = {bf},
  font            = {normalsize},     % 10pt
  format          = hang,             % baris ke-2 sejajar judul (6c)
  justification   = raggedright,
  singlelinecheck = false,
  margin          = 1cm,              % menjorok 1 cm kiri-kanan (6b)
  skip            = 12pt,             % jarak 1 spasi (6i)
  name            = Gambar,
  labelsep        = colon             % "Gambar 2.1:"
}

% -- Caption Tabel (Pasal 4.2.2 no. 5) ------------------------
% "Tabel 2.1:" bold, judul 10pt normal, di atas tabel, menjorok
% 1 cm dari margin kiri-kanan, baris lanjutan sejajar judul.
\captionsetup[table]{
  labelfont       = {bf},             % "Tabel 2.1:" bold (5b)
  textfont        = {},               % judul 10pt normal (5e)
  font            = {normalsize},
  format          = hang,             % baris ke-2 sejajar judul (5c)
  justification   = raggedright,
  singlelinecheck = false,
  margin          = 1cm,              % menjorok 1 cm kiri-kanan (5b)
  skip            = 12pt,             % jarak 1 spasi (5i)
  position        = above,
  name            = Tabel,
  labelsep        = colon             % "Tabel 2.1:"
}
\renewcommand{\tablename}{Tabel}

% -- Format Default Tabularray --------------------------------
\SetTblrInner{
  hlines, vlines,
  row{1} = {bg=gray!20, font=\bfseries, halign=c},
  rowhead = 1,
}

\SetTblrInner[longtblr]{
  hlines, vlines,
  row{1} = {bg=gray!20, font=\bfseries, halign=c},
  rowhead = 1,
}

% -- Isi Tabel: lebih kecil (\small, +-9pt), spasi 1 ----------
% Berlaku otomatis untuk semua tblr/longtblr. Caption tidak terpengaruh.
\AddToHook{env/tblr/begin}{\small\linespread{1}\selectfont}
\AddToHook{env/longtblr/begin}{\small\linespread{1}\selectfont}

% -- Kustomisasi Khusus longtblr ------------------------------
% Caption longtblr disamakan dengan tabel biasa: "Tabel 2.1:" bold,
% judul 10pt normal, rata kiri.
\DefTblrTemplate{caption-sep}{default}{:\space}
\SetTblrStyle{caption-tag}{font=\normalsize\bfseries}
\SetTblrStyle{caption-text}{font=\normalsize}
\SetTblrStyle{caption-halign}{l}

\DefTblrTemplate{contfoot-text}{default}{}
\DefTblrTemplate{conthead-text}{default}{}
\DefTblrTemplate{capcont}{default}{}
\DefTblrTemplate{contfoot}{default}{}
\DefTblrTemplate{conthead}{default}{}

% -- Penomoran Gambar/Tabel/Persamaan per Bab -----------------
\renewcommand{\thefigure}{\arabic{chapter}.\arabic{figure}}
\renewcommand{\thetable}{\arabic{chapter}.\arabic{table}}
\renewcommand{\theequation}{\arabic{chapter}.\arabic{equation}}

% -- Header & Footer ------------------------------------------
\setlength{\headheight}{40pt}  % ruang untuk logo di header
\pagestyle{fancy}
\fancyhf{}
\fancyfoot[C]{\thepage}
\renewcommand{\headrulewidth}{0pt}


% Logo PENS untuk header Gaya 2 (diukur sekali agar lebar teks di
% sebelahnya bisa dihitung)
\newsavebox{\KotakLogoHeader}
\sbox{\KotakLogoHeader}{\includegraphics[height=1.0cm]{assets/pens/logo_pens.png}}

% Pagestyle kpruler dengan 2 macam gaya (Simple & PENS)
\fancypagestyle{kpruler}{
  \fancyhf{}
  \ifnum\StyleHeaderFooter=2
    % Gaya 2: PENS (Gambar 2)
    % Nama laporan & perusahaan dipindah ke header (baris ke-2) agar
    % nomor halaman bisa di bawah tengah (Pasal 4.2.2 no. 4c).
    % Logo di kiri, teks di ruang SISA sebelah kanan logo (rata tengah
    % di ruang itu) agar teks tidak menimpa logo. Baris yang terlalu
    % panjang otomatis dikecilkan (\fitline). Nama perusahaan di header memakai
    % \NamaPerusahaanSingkat (config/variables.tex) agar ringkas.
    \fancyhead[LO,LE]{%
      \raisebox{-0.2cm}{\usebox{\KotakLogoHeader}}\hspace{2mm}%
      \parbox[b]{\dimexpr\headwidth-\wd\KotakLogoHeader-2mm\relax}{%
        \centering\FontHeaderFooterPENS
        \fitline{\linewidth}{POLITEKNIK ELEKTRONIKA NEGERI SURABAYA}\\
        \fitline{\linewidth}{LAPORAN KERJA PRAKTIK~-~\MakeUppercase{\ifcsvoid{NamaPerusahaanSingkat}{\NamaPerusahaan}{\NamaPerusahaanSingkat}}}%
        \rule[-3mm]{0pt}{0pt}}}
    \renewcommand{\headrulewidth}{1.5pt}
    \fancyfoot[C]{\FontHeaderFooterPENS\thepage}
    \renewcommand{\footrulewidth}{1.5pt}
  \else
    % Gaya 1: Simple (Gambar 1)
    \fancyhead[C]{\FontHeaderFooterSimple\bfseries LAPORAN KERJA PRAKTIK\rule[-3mm]{0pt}{0pt}}
    \renewcommand{\headrulewidth}{1.5pt}
    \fancyfoot[C]{\FontHeaderFooterSimple\thepage}
    \renewcommand{\footrulewidth}{1.5pt}
  \fi
}

% Redefine plain to copy kpruler for consistency on chapter start pages
\makeatletter
\let\ps@plain\ps@kpruler
\makeatother

% -- Lingkungan Kutipan Panjang (lebih dari 2 baris) -----------
\newenvironment{kutipanpanjang}{%
  \begin{list}{}{%
      \setlength{\leftmargin}{10mm}
      \setlength{\rightmargin}{10mm}
      \setlength{\topsep}{6pt}
      \setlength{\itemsep}{0pt}
      \setlength{\parsep}{0pt}
    }
    \item\FontKutipanPanjang\ignorespaces
    }{%
    \unskip
  \end{list}
}

% -- Hyperref -------------------------------------------------
\hypersetup{
  colorlinks  = true,
  linkcolor   = black,
  citecolor   = black,
  urlcolor    = black,
  plainpages  = false,   % cegah duplicate anchor saat roman->arabic reset
  pdftitle    = {\JudulKP},
  pdfauthor   = {\NamaMahasiswaSatu}
}

% -- Nama Bahasa Indonesia ------------------------------------
\renewcommand{\contentsname}{DAFTAR ISI}
\renewcommand{\listfigurename}{DAFTAR GAMBAR}
\renewcommand{\listtablename}{DAFTAR TABEL}
\renewcommand{\bibname}{DAFTAR PUSTAKA}
\renewcommand{\appendixname}{LAMPIRAN}
\renewcommand{\chaptername}{BAB}

% -- Judul Daftar Isi/Gambar/Tabel: 12pt, Bold, Center --------
\renewcommand{\cfttoctitlefont}{\hfill\FontJudulDaftar\bfseries\MakeUppercase}
\renewcommand{\cftaftertoctitle}{\hfill\null}

\renewcommand{\cftloftitlefont}{\hfill\FontJudulDaftar\bfseries\MakeUppercase}
\renewcommand{\cftafterloftitle}{\hfill\null}

\renewcommand{\cftlottitlefont}{\hfill\FontJudulDaftar\bfseries\MakeUppercase}
\renewcommand{\cftafterlottitle}{\hfill\null}

\setlength{\cftbeforetoctitleskip}{0pt}
\setlength{\cftbeforeloftitleskip}{0pt}
\setlength{\cftbeforelottitleskip}{0pt}
\setlength{\cftaftertoctitleskip}{10pt}
\setlength{\cftafterloftitleskip}{10pt}
\setlength{\cftafterlottitleskip}{10pt}

% -- Prefix "Gambar" dan "Tabel" di LOF/LOT -------------------
\renewcommand{\cftfigpresnum}{Gambar }
\renewcommand{\cfttabpresnum}{Tabel }
% numwidth harus cukup lebar untuk menampung prefix terpanjang
% agar baris kedua (wrap) sejajar dengan teks, bukan kembali ke kiri.
% "Gambar 4.12 " ≈ 7em; "Tabel 4.13 " ≈ 6em
\setlength{\cftfignumwidth}{7.5em}
\setlength{\cfttabnumwidth}{6.5em}

% -- Titik-titik di Daftar Isi --------------------------------
\renewcommand{\cftdotsep}{1.5}

% -- Spasi Konsisten di Daftar Isi ----------------------------
\setlength{\cftbeforechapskip}{0pt}
\setlength{\cftbeforesecskip}{0pt}
\setlength{\cftbeforesubsecskip}{0pt}
\setlength{\cftbeforefigskip}{0pt}
\setlength{\cftbeforetabskip}{0pt}

% -- Font Level Bab: Bold -------------------------------------
\renewcommand{\cftchapfont}{\bfseries}
\renewcommand{\cftchappagefont}{\bfseries}
\renewcommand{\cftchapleader}{\cftdotfill{\cftdotsep}}

% -- Font Level Section: Normal -------------------------------
\renewcommand{\cftsecfont}{\normalfont}
\renewcommand{\cftsecpagefont}{\normalfont}

% -- Prefix "BAB" di Daftar Isi -------------------------------
\renewcommand{\cftchappresnum}{BAB }
% "BAB IV " ≈ 5em -- set numwidth eksplisit agar wrap sejajar teks
\setlength{\cftchapnumwidth}{5.5em}
% Indent section/subsection dikembalikan ke default agar menjorok bertingkat
% \setlength{\cftsecindent}{1.5em}
% \setlength{\cftsubsecindent}{3.8em}

% -- Jarak (Spasi) Atas/Bawah Gambar dan Tabel ----------------
\setlength{\intextsep}{10pt}
\setlength{\textfloatsep}{10pt}

\setlength{\bibsep}{3pt}

% -- Baris Yatim/Janda ----------------------------------------
% (Hyphenation sudah dinonaktifkan di bagian atas file ini.)
% Cegah satu baris paragraf tertinggal sendirian di bawah halaman
% (club/yatim) atau di atas halaman berikutnya (widow/janda).
\clubpenalty=10000
\widowpenalty=10000
\displaywidowpenalty=10000

% -- Halaman Kosong -------------------------------------------
% Halaman kosong pada cetak bolak-balik diberi keterangan
\makeatletter
\newif\if@mainmatter
\@mainmatterfalse
\def\cleardoublepage{%
  \clearpage
  \if@twoside
    \ifodd\c@page\else
      \hbox{}
      \vspace*{\fill}
      \begin{center}
        \textit{Halaman ini sengaja dikosongkan}
      \end{center}
      \vspace{\fill}
      \thispagestyle{kpruler}
      \newpage
      \if@twocolumn\hbox{}\newpage\fi
    \fi
  \fi
}

% -- Paksa Halaman Ganjil -------------------------------------
\newcommand{\cleartooddpage}{%
  \clearpage
  \ifodd\c@page\else
    \hbox{}
    \vspace*{\fill}
    \begin{center}
      \textit{Halaman ini sengaja dikosongkan}
    \end{center}
    \vspace{\fill}
    \thispagestyle{fancy}
    \newpage
  \fi
}
\makeatother

%============================================================
% FORMAT: BUKU KERJA PRAKTIK
% Aturan layout khusus Buku KP PENS.
%
% Kertas  : A5 (148 mm x 210 mm)
% Font    : 10pt normal, BAB 12pt (diset di formatting.tex)
% Spasi   : 1 spasi (single)
% Cetak   : Bolak-balik (twoside), halaman ganjil = depan
% Margin  : Mirror -- sisi jilid 2.5 cm, sisi luar 2 cm
%   Atas  : 2.5 cm
%   Bawah : 2.0 cm
%   Inner : 2.5 cm (sisi jilid -- kiri halaman ganjil)
%   Outer : 2.0 cm (sisi luar)
%============================================================

% -- Geometri: A5, double-sided, mirror margin -----------------
\geometry{
  a5paper,
  twoside,
  top    = 2.5cm,
  bottom = 2.0cm,
  inner  = 2.5cm,   % Margin sisi jilid (kiri halaman ganjil)
  outer  = 2.0cm    % Margin luar (kanan halaman ganjil)
}

% -- Spasi Baris: 1 spasi (single) ----------------------------
\singlespacing

% -- Indentasi Paragraf ----------------------------------------
\setlength{\parindent}{1.0cm}


%============================================================
\begin{document}
%============================================================

\hypersetup{pageanchor=false}

% -- Halaman i: Cover (romawi i, tanpa header/footer) -----
% PENTING: kpruler diset GLOBAL sebelum cover.
% \begin{titlepage} di cover_buku.tex otomatis override dengan \thispagestyle{empty}
% sehingga cover tidak punya header/footer.
% Setelah \end{titlepage}, global pagestyle tetap kpruler.
\pagenumbering{roman}
\setcounter{page}{1}
\pagestyle{kpruler}
%============================================================
% HALAMAN SAMPUL BUKU KP (Cover Berwarna)
% Layout mengikuti Template_Buku_PA_PENS_LaTeX:
%   - Header putih: logo PENS (kiri) + "KERJA PRAKTIK" (kanan, orange)
%   - Body orange: judul kapital bold, Oleh:, mahasiswa, pembimbing, footer
% Mendukung 1, 2, atau 3 mahasiswa.
%============================================================
\begin{titlepage}
  \newgeometry{top=0pt, bottom=0pt, left=0pt, right=0pt}
  \setlength{\fboxsep}{0pt}

  % -- Bagian 1: Header Putih ----------------------------------------
  \noindent\pagecolor{white}
  \colorbox{white}{%
    \begin{minipage}[c][5.36cm][c]{\paperwidth}
      \hspace*{1.2cm}
      \begin{minipage}[t]{2.8cm}
        \vspace{0pt} % Set baseline at top
        \raisebox{-\height}{\includegraphics[width=2.5cm, keepaspectratio]{assets/pens/logo_pens.png}}
      \end{minipage}
      \hfill
      \begin{minipage}[t]{5cm}
        \vspace{0pt} % Set baseline at top
        \raggedleft
        \vspace{0.25cm}
        \raisebox{-\height}{\color{black}\FontCoverLabel\bfseries KERJA PRAKTIK}
      \end{minipage}
      \hspace*{1.2cm}
    \end{minipage}
  }\par\nointerlineskip

  % -- Bagian 2: Body Orange (latar belakang penuh) ------------------
  \noindent
  \colorbox{penskuning}{%
    \begin{minipage}[t][\dimexpr\paperheight-5.36cm\relax][t]{\paperwidth}
      \color{black}
      \begin{center}

        \vspace{0.8cm}

        % Judul: kapital, bold (TANPA italic), kata tidak dipotong
        \begin{minipage}{\LebarJudulKP}
          \centering
          \hyphenpenalty=10000
          \exhyphenpenalty=10000
          \FontCoverJudul\bfseries
          \MakeUppercase{\JudulKP}
        \end{minipage}
        \par

        \vfill

        % Blok Penulis
        {\FontCoverIsi\itshape Oleh:\par}
        \vspace{0.2cm}

        \begin{tabular}{@{}p{5.5cm}l@{}}
          \bfseries\FontCoverIsi \NamaMahasiswaSatu & \bfseries\FontCoverIsi NRP. \NRPSatu \\
          \ifcsvoid{NamaMahasiswaDua}{}{%
          \bfseries\FontCoverIsi \NamaMahasiswaDua  & \bfseries\FontCoverIsi NRP. \NRPDua  \\
          }
          \ifcsvoid{NamaMahasiswaTiga}{}{%
          \bfseries\FontCoverIsi \NamaMahasiswaTiga & \bfseries\FontCoverIsi NRP. \NRPTiga \\
          }
        \end{tabular}

        \vfill

        % Dosen Pembimbing
        {\FontCoverIsi\itshape Dosen Pembimbing:\par}
        \vspace{0.15cm}
        {\FontCoverIsi\bfseries \NamaPembimbing\par}
        {\FontCoverIsi\bfseries NIP. \NIPPembimbing\par}

        \vfill

        % Footer: institusi
        {\FontCoverIsi\bfseries \MakeUppercase{\JenjangProdi} \MakeUppercase{\NamaProdi}\par}
        {\FontCoverIsi\bfseries \MakeUppercase{\NamaDepartemen}\par}
        {\FontCoverIsi\bfseries POLITEKNIK ELEKTRONIKA NEGERI SURABAYA\par}
        {\FontCoverIsi\bfseries \TahunKP\par}
        \vspace{1.5cm}

      \end{center}
    \end{minipage}%
  }
\end{titlepage}
\restoregeometry
\pagecolor{white}


% -- Halaman ii: Blank, dibuat otomatis oleh \cleardoublepage -----
% (pakai mekanisme yg SAMA persis dgn hal. iv yang sudah terbukti benar)
\makeatletter\@mainmatterfalse\makeatother
\cleardoublepage

% -- Halaman iii: Halaman Judul Dalam (romawi iii, tanpa header/footer)
%============================================================
% HALAMAN JUDUL DALAM BUKU KP
% Header/footer ditangani oleh pagestyle kpruler.
% Mendukung 1, 2, atau 3 mahasiswa.
%============================================================
\thispagestyle{empty}

\begin{center}
  \vspace{0.3cm}

  {\FontJudulDalamLabel\bfseries LAPORAN KERJA PRAKTIK\par}
  \vspace{1cm}

  \begin{minipage}{0.92\textwidth}
    \centering
    {\FontJudulDalamJudul\bfseries\MakeUppercase{\JudulKP}\par}
  \end{minipage}
  
  \vspace{0.8cm}
  {\FontJudulDalamIsi\bfseries Periode \PeriodeKP\par}

  \vfill

  \includegraphics[height=5.5cm]{assets/pens/logo_pens.png}\par

  \vfill

  {\FontJudulDalamIsi\bfseries Oleh:\par}
  \vspace{0.3cm}

  \begin{tabular}{@{}p{5.5cm}l@{}}
    \bfseries\FontJudulDalamIsi \NamaMahasiswaSatu & \bfseries\FontJudulDalamIsi NRP. \NRPSatu \\
    \ifcsvoid{NamaMahasiswaDua}{}{%
      \bfseries\FontJudulDalamIsi \NamaMahasiswaDua & \bfseries\FontJudulDalamIsi NRP. \NRPDua \\
    }
    \ifcsvoid{NamaMahasiswaTiga}{}{%
      \bfseries\FontJudulDalamIsi \NamaMahasiswaTiga & \bfseries\FontJudulDalamIsi NRP. \NRPTiga \\
    }
  \end{tabular}

  \vfill

  {\FontJudulDalamIsi\bfseries \MakeUppercase{\JenjangProdi} \MakeUppercase{\NamaProdi}\par}
  {\FontJudulDalamIsi\bfseries \MakeUppercase{\NamaDepartemen}\par}
  {\FontJudulDalamIsi\bfseries POLITEKNIK ELEKTRONIKA NEGERI SURABAYA\par}
  {\FontJudulDalamIsi\bfseries \TahunKP\par}

\end{center}
\thispagestyle{empty} % tanpa nomor halaman (Pasal 4.2.2 no. 4b)
  % TIDAK ada \newpage lagi di dalam file ini

% -- Halaman iii: Lembar Pengesahan Internal ------------------
% (halaman ii = blank dengan kpruler)
\cleardoublepage
\ifpengesahaninternalpdf
  \includepdf[pages=-]{\FilePengesahanInternalPDF}
\else
  %============================================================
% LEMBAR PENGESAHAN INTERNAL (PENS)
% Ditandatangani: Dosen Pembimbing, Koordinator KP (sejajar)
% Di bawahnya: Kepala Prodi D4 Teknik Informatika
%============================================================
\newpage

\begin{center}
  {\FontPengesahanJudul\bfseries LEMBAR PENGESAHAN\par}
  \vspace{0.3cm}
  {\FontPengesahanJudul\bfseries LAPORAN KERJA PRAKTIK\par}
  {\FontPengesahanJudul\bfseries \MakeUppercase{\NamaPerusahaan}\par}
  {\FontPengesahanJudul Periode \PeriodeKP\par}
  \vspace{1cm}

  {\FontPengesahanIsi Disusun oleh:\par}
  \vspace{0.1cm}

  \begin{tabular}{@{}p{5.5cm}l@{}}
    \bfseries\FontPengesahanIsi \NamaMahasiswaSatu & \bfseries\FontPengesahanIsi \NRPSatu \\
    \ifcsvoid{NamaMahasiswaDua}{}{%
      \bfseries\FontPengesahanIsi \NamaMahasiswaDua & \bfseries\FontPengesahanIsi \NRPDua \\
    }
    \ifcsvoid{NamaMahasiswaTiga}{}{%
      \bfseries\FontPengesahanIsi \NamaMahasiswaTiga & \bfseries\FontPengesahanIsi \NRPTiga \\
    }
  \end{tabular}
\end{center}

\vspace{2cm}

\begin{center}{\FontPengesahanIsi\bfseries Menyetujui,}\end{center}
\vspace{0.2cm}

\noindent
\begin{minipage}[t]{0.46\textwidth}
  \centering
  {\FontPengesahanIsi Dosen Pembimbing\par}
  \vspace{1.6cm}
  {\bfseries\FontPengesahanIsi \underline{\NamaPembimbing}}\par
  \vspace{0.15cm}
  {\FontPengesahanIsi NIP. \NIPPembimbing\par}
\end{minipage}%
\hfill
\begin{minipage}[t]{0.46\textwidth}
  \centering
  {\FontPengesahanIsi Koordinator Kerja Praktik\par}
  \vspace{1.6cm}
  {\bfseries\FontPengesahanIsi \underline{\NamaKoordinatorKP}}\par
  \vspace{0.15cm}
  {\FontPengesahanIsi NIP. \NIPKoordinatorKP\par}
\end{minipage}

\vspace{2cm}

\begin{center}{\FontPengesahanIsi\bfseries Mengetahui,}\end{center}
\vspace{0.2cm}

\begin{center}
  {\FontPengesahanIsi Kepala Program Studi D4 Teknik Informatika\par}
  \vspace{1.6cm}
  {\bfseries\FontPengesahanIsi \underline{\NamaKetuaProdi}}\par
  \vspace{0.15cm}
  {\FontPengesahanIsi NIP. \NIPKetuaProdi\par}
\end{center}
\thispagestyle{empty} % tanpa nomor halaman (Pasal 4.2.2 no. 4b)

\fi

% -- Halaman vii: Lembar Pengesahan Instansi ------------------
% (halaman vi = blank)
\cleardoublepage
\ifpengesahaninstansipdf
  \includepdf[pages=-]{\FilePengesahanInstansiPDF}
\else
  %============================================================
% LEMBAR PENGESAHAN INSTANSI/PERUSAHAAN
% Ditandatangani: Pembimbing Lapangan dari perusahaan
%============================================================
\newpage

\begin{center}
  {\FontPengesahanJudul\bfseries LEMBAR PENGESAHAN\par}
  \vspace{0.3cm}
  {\FontPengesahanJudul\bfseries LAPORAN KERJA PRAKTIK\par}
  {\FontPengesahanJudul\bfseries \MakeUppercase{\NamaPerusahaan}\par}
  {\FontPengesahanJudul Periode \PeriodeKP\par}
  \vspace{1cm}

  {\FontPengesahanIsi Disusun oleh:\par}
  \vspace{0.1cm}

  \begin{tabular}{@{}p{5.5cm}l@{}}
    \bfseries\FontPengesahanIsi \NamaMahasiswaSatu & \bfseries\FontPengesahanIsi \NRPSatu \\
    \ifcsvoid{NamaMahasiswaDua}{}{%
      \bfseries\FontPengesahanIsi \NamaMahasiswaDua & \bfseries\FontPengesahanIsi \NRPDua \\
    }
    \ifcsvoid{NamaMahasiswaTiga}{}{%
      \bfseries\FontPengesahanIsi \NamaMahasiswaTiga & \bfseries\FontPengesahanIsi \NRPTiga \\
    }
  \end{tabular}
\end{center}

\vspace{2.5cm}

\begin{center}
  {\FontPengesahanIsi\bfseries Menyetujui,\par}
  \vspace{0.3cm}

  {\FontPengesahanIsi PEMBIMBING KERJA PRAKTIK\par}
  {\FontPengesahanIsi \MakeUppercase{\JabatanPembimbingLapangan}\par}
  \vspace{2.2cm}
  {\FontPengesahanIsi\bfseries \underline{\NamaPembimbingLapangan}\par}
\end{center}
\thispagestyle{empty} % tanpa nomor halaman (Pasal 4.2.2 no. 4b)

\fi

% -- Halaman ix: Abstrak ---------------------------------------
% (halaman viii = blank)
\cleardoublepage
\makeatletter\@mainmattertrue\makeatother
%============================================================
% HALAMAN ABSTRAK  (ringkasan dalam Bahasa Indonesia)
% Pedoman KP Pasal 4.2.1 (iv): 150-200 kata, maksimal satu
% halaman, 1,5 spasi. Isi: apa yang dikerjakan, metodologi,
% hasil yang dicapai; jika menyangkut program, sebutkan
% platform dan development tool yang dipakai.
%============================================================
\cleardoublepage
\phantomsection
\addcontentsline{toc}{chapter}{ABSTRAK}
\begin{center}
  {\FontJudulHalaman\bfseries ABSTRAK\par}
\end{center}
\vspace{12pt}
% Spasi 1,5 memakai \setstretch{1.5} (lewat lingkungan spacing),
% sama seperti Buku PA. \onehalfspacing / onehalfspace jauh lebih
% rapat dari spasi 1,5 MS Word, jadi tidak dipakai.
\begin{spacing}{1.5}

Laporan ini membahas pelaksanaan Kerja Praktik di PT. Sinergi Dimensi Informatika
(SINDIKA) yang berfokus pada pengembangan sistem monitoring infrastruktur kelistrikan
`Voltra' dan library backend Nedo.AspNet.Request.Validation. Pada sistem Voltra, penulis
mengimplementasikan kalkulasi Health Index, Critical Index, dan Risk Index untuk
klasifikasi status aset, manajemen tiket insiden dengan auto-close berbasis event, analitik
finansial (cost avoidance dan risk cost projection), evaluasi keusangan aset beserta
prediksi Remaining Useful Life (RUL), serta pelaporan otomatis berbasis PDF. Sistem
dibangun dengan arsitektur microservices menggunakan RabbitMQ sebagai message broker, MQTT
untuk ingesti data sensor, Redis sebagai cache dan pub/sub, serta InfluxDB dan PostgreSQL
untuk data time-series dan relasional. Penulis juga mengembangkan
Nedo.AspNet.Request.Validation, library validasi input berbasis attribute untuk ASP.NET
Core yang menyediakan validasi terpusat dan konsisten, standarisasi kode error, serta
dukungan multi-bahasa pada respons API. Pengembangan dilakukan secara iteratif bersama tim
melalui alur kerja Git, code review, dan pengujian fungsional. Kegiatan ini memberikan pemahaman mendalam mengenai
penerapan arsitektur microservices, komunikasi asinkron, dan pengembangan library internal
dalam skenario industri nyata.

\end{spacing}

\vspace{0.5cm}
\noindent\textbf{Kata Kunci}

\noindent Voltra, Microservices, RabbitMQ, MQTT, Health Index, Risk Index, Library Validation, ASP.NET Core, Predictive Maintenance


% -- Urutan Pasal 4.2.1: Daftar Isi, Kata Pengantar, Daftar ---
% -- Gambar, Daftar Tabel, Daftar Lampiran ---------------------
\cleardoublepage
\addcontentsline{toc}{chapter}{DAFTAR ISI}
\tableofcontents

\cleardoublepage
%============================================================
% HALAMAN KATA PENGANTAR
% Gunakan bahasa formal. Sertakan ucapan terima kasih kepada
% pihak-pihak yang berkontribusi pada kerja praktik.
%============================================================
\cleardoublepage
\phantomsection
\addcontentsline{toc}{chapter}{KATA PENGANTAR}
\begin{center}
  {\FontJudulHalaman\bfseries KATA PENGANTAR\par}
\end{center}
\vspace{12pt}

Puji Syukur kami ucapkan kepada Tuhan Yang Maha Esa atas limpahan berkat dan rahmat-Nya
sehingga kami dapat menyelesaikan Laporan Kerja Praktik di \NamaPerusahaan\ yang
dilaksanakan mulai tanggal \PeriodeKP\ dengan judul \textit{``\JudulKP''}.

Laporan ini merupakan salah satu syarat bagi penyelesaian kerja praktik jurusan Teknik
Informatika PENS (Politeknik Elektronika Negeri Surabaya), dan juga sebagai sarana
pembelajaran dan pendalaman untuk menyesuaikan ilmu pengetahuan teoritis yang diterima di
perkuliahan dengan pengetahuan nyata di dunia kerja.

Selama proses pembuatan laporan kerja praktik, banyak bantuan, saran, dan motivasi yang
diberikan oleh berbagai pihak sehingga kami dapat menyelesaikan laporan ini. Oleh karena
itu, kami mengucapkan terima kasih kepada:

\begin{enumerate}[label=\arabic*., leftmargin=*, itemsep=0pt]
  \item Allah SWT, Tuhan Yang Maha Esa yang telah melimpahkan kemudahan kepada penulis
        sejak dalam pelaksanaan tugas kerja praktik di \NamaPerusahaan\ sampai dengan
        penyusunan laporan kerja praktik.
  \item Orangtua dan keluarga yang selalu memberikan dukungan moril maupun materiil sejak
        kerja praktik hingga penyusunan laporan kerja praktik.
  \item Bapak M. Udin Harun Al Rasyid, S.Kom., Ph.D. selaku Kepala Departemen Teknik
        Informatika dan Komputer PENS.
  \item Bapak \NamaKetuaProdi\ selaku Ketua Prodi \JenjangProdi\ \NamaProdi\ PENS.
  \item Ibu \NamaKoordinatorKP\ selaku Koordinator Kerja Praktik Jurusan Teknik Informatika PENS.
  \item Ibu \NamaPembimbing\ selaku Dosen Pembimbing yang telah meluangkan waktu dan
        memberi bimbingan, saran, dan arahan dalam proses pengerjaan proyek maupun saat KP
        berlangsung.
  \item \NamaPerusahaan\ yang telah memberikan kesempatan pada penulis untuk melaksanakan
        kerja praktik.
  \item Bapak \NamaPembimbingLapangan\ selaku \JabatanPembimbingLapangan\ sekaligus
        Pembimbing Lapangan di \NamaPerusahaan.
  \item Para karyawan/karyawati di \NamaPerusahaan\ yang telah bersedia bekerja sama dan
        memberikan informasi-informasi yang dibutuhkan dalam pelaksanaan kerja praktik ini.
  \item Serta pihak-pihak lain yang telah membantu kami namun tidak dapat disebutkan satu
        persatu.
\end{enumerate}

Demikian laporan kerja praktik ini disusun. Kami menyadari masih banyak kekurangan dalam
penyusunan laporan ini. Kami berharap semoga apa yang kami kerjakan dan kami tulis dalam
laporan ini dapat membantu memberi manfaat bagi pihak-pihak terkait.

\vspace{1cm}
\begin{flushright}
  \TanggalPengesahan\\[0.8cm]
  Penulis
\end{flushright}


\cleardoublepage
\begingroup
\renewcommand*{\addvspace}[1]{}
\listoffigures
\endgroup
\addcontentsline{toc}{chapter}{DAFTAR GAMBAR}

\cleardoublepage
\begingroup
\renewcommand*{\addvspace}[1]{}
\listoftables
\endgroup
\addcontentsline{toc}{chapter}{DAFTAR TABEL}

%============================================================
% DAFTAR LAMPIRAN (Pasal 4.2.1 ix)
% Urutan harus sama dengan \section di pages/lampiran.tex.
%============================================================
\cleardoublepage
\phantomsection
\addcontentsline{toc}{chapter}{DAFTAR LAMPIRAN}
\begin{center}
  {\FontJudulHalaman\bfseries DAFTAR LAMPIRAN\par}
\end{center}
\vspace{12pt}

\noindent
\begin{tabularx}{\linewidth}{@{}l X@{}}
  Lampiran A & Rekap Monitoring Kegiatan KP yang Telah Diverifikasi \\
  Lampiran B & Nilai Pembimbing Perusahaan \\
  Lampiran C & Nilai Dosen Pembimbing KP \\
  Lampiran D & Nilai Dosen Penguji KP \\
  Lampiran E & Foto Mahasiswa KP di Lokasi KP \\
  Lampiran F & Biodata Penulis \\
\end{tabularx}


% -- Bab 1+: Arabic dari 1, pageanchor aktif ------------------
\cleardoublepage
\pagenumbering{arabic}
\setcounter{page}{1}
\hypersetup{pageanchor=true}
\pagestyle{kpruler}

\newcommand{\SistematikaFile}{pages/sistematika_kp}
%============================================================
% BAB 1: PENDAHULUAN
%   1.1 Latar Belakang
%   1.2 Rumusan Masalah
%   1.3 Batasan Masalah
%   1.4 Tujuan
%       1.4.1 Tujuan Proyek [Nama Proyek 1]
%       1.4.2 Tujuan Proyek [Nama Proyek 2] (hapus jika hanya 1 proyek)
%   1.5 Manfaat
%       1.5.1 Manfaat Proyek [Nama Proyek 1]
%       1.5.2 Manfaat Proyek [Nama Proyek 2] (hapus jika hanya 1 proyek)
%   1.6 Sistematika Penulisan
%============================================================
\chapter{PENDAHULUAN}

\section{Latar Belakang}

Deskripsikan latar belakang pemilihan tempat dan topik kerja praktik. Jelaskan
konteks industri atau institusi tempat KP dilaksanakan, permasalahan yang ada,
relevansi dengan program studi, serta gambaran umum proyek yang dikerjakan.
Latar belakang yang baik mampu menjelaskan urgensi dan motivasi pelaksanaan
kerja praktik ini.

\section{Rumusan Masalah}

Berdasarkan latar belakang yang telah dijelaskan, maka permasalahan yang akan
dibahas dalam laporan kerja praktik ini dapat dirumuskan sebagai berikut.

\begin{enumerate}[label=\alph*., leftmargin=*, itemsep=0pt]
  \item Rumusan masalah pertama.
  \item Rumusan masalah kedua.
  \item Rumusan masalah ketiga.
\end{enumerate}

\section{Batasan Masalah}

Untuk memfokuskan pembahasan pada laporan kerja praktik ini, maka ditetapkan
batasan masalah sebagai berikut:

\begin{enumerate}[label=\alph*., leftmargin=*, itemsep=0pt]
  \item Pembahasan difokuskan pada kegiatan yang dilakukan selama periode
        \PeriodeKP\ di \NamaPerusahaan.
  \item Batasan teknis kedua (misalnya: hanya mencakup sisi frontend/backend).
  \item Batasan ketiga (aspek yang tidak dibahas dalam laporan ini).
  \item Periode pengamatan dan keterlibatan dalam proyek dibatasi selama masa
        kerja praktik.
\end{enumerate}

\section{Tujuan}

Bagian ini menjelaskan tujuan dari pelaksanaan proyek kerja praktik yang
dilakukan di \NamaPerusahaan. Tujuan yang dirumuskan tidak hanya menggambarkan
hasil yang ingin dicapai dari sisi teknis, tetapi juga mencerminkan kontribusi
sistem terhadap efisiensi kerja dan nilai praktis dalam dunia industri.

% Jika hanya satu proyek, hapus sub-section 1.4.2 dan ubah 1.4.1 jadi narasi langsung.
\subsection{Tujuan Proyek [Nama Proyek 1]}

\begin{itemize}[leftmargin=*, itemsep=0pt]
  \item Tujuan pertama proyek.
  \item Tujuan kedua proyek.
  \item Tujuan ketiga proyek.
\end{itemize}

\subsection{Tujuan Proyek [Nama Proyek 2]}

\begin{itemize}[leftmargin=*, itemsep=0pt]
  \item Tujuan pertama proyek kedua.
  \item Tujuan kedua proyek kedua.
\end{itemize}

\section{Manfaat}

Pelaksanaan kerja praktik di \NamaPerusahaan\ tidak hanya bertujuan untuk
memenuhi kewajiban akademik, tetapi juga memberikan manfaat nyata baik dari
sisi teknis, operasional, maupun strategis. Berikut adalah manfaat utama dari
masing-masing proyek:

% Jika hanya satu proyek, hapus sub-section 1.5.2 dan ubah 1.5.1 jadi narasi langsung.
\subsection{Manfaat Proyek [Nama Proyek 1]}

\begin{itemize}[leftmargin=*, itemsep=0pt]
  \item \textbf{Manfaat Pertama.} Deskripsi manfaat pertama.
  \item \textbf{Manfaat Kedua.} Deskripsi manfaat kedua.
  \item \textbf{Manfaat Ketiga.} Deskripsi manfaat ketiga.
\end{itemize}

\subsection{Manfaat Proyek [Nama Proyek 2]}

\begin{itemize}[leftmargin=*, itemsep=0pt]
  \item \textbf{Manfaat Pertama.} Deskripsi manfaat pertama proyek kedua.
  \item \textbf{Manfaat Kedua.} Deskripsi manfaat kedua proyek kedua.
\end{itemize}

\section{Sistematika Penulisan}

Berikut sistematika penulisan pada laporan kerja praktik ini:

\input{\SistematikaFile}

%============================================================
% BAB 2: GAMBARAN PERUSAHAAN
%   2.1 Profil Perusahaan (termasuk visi & misi dalam narasi)
%   2.2 Struktur Organisasi Perusahaan
%       2.2.x per divisi
%   2.3 Hak dan Wewenang
%       2.3.x per jabatan
%   2.4 Lokasi Perusahaan
%   2.5 Etika Profesi
%============================================================
\chapter{GAMBARAN PERUSAHAAN}

\section{Profil Perusahaan}

Jelaskan profil singkat \NamaPerusahaan, meliputi: tahun berdiri, bidang usaha
utama, skala perusahaan, dan gambaran umum posisinya di industri atau sektor
yang bersangkutan.

Sertakan juga visi dan misi perusahaan dalam bagian ini sebagai narasi
atau daftar ringkas. Jelaskan nilai-nilai utama yang dijalankan perusahaan.

% Contoh penyisipan logo:
% \begin{figure}[H]
%   \centering
%   \includegraphics[width=4cm]{assets/gambar/logo_perusahaan.png}
%   \caption{Logo \NamaPerusahaan}
%   \label{fig:logo_perusahaan}
% \end{figure}

\section{Struktur Organisasi Perusahaan}

Struktur organisasi \NamaPerusahaan\ dirancang untuk mendukung efisiensi
operasional dan kolaborasi antar departemen. Berikut adalah struktur organisasi
\NamaPerusahaan.

% Contoh penyisipan bagan organisasi:
% \begin{figure}[H]
%   \centering
%   \includegraphics[width=0.9\textwidth]{assets/gambar/struktur_organisasi.png}
%   \caption{Struktur Organisasi \NamaPerusahaan}
%   \label{fig:struktur_org}
% \end{figure}

% Sesuaikan sub-section dengan divisi yang ada di perusahaan.
\subsection{[Nama Divisi 1]}

\begin{itemize}[leftmargin=*, itemsep=0pt]
  \item \textbf{Kepala Divisi:} Nama kepala divisi.
  \item \textbf{Tugas:} Deskripsi tugas dan tanggung jawab divisi ini.
\end{itemize}

\subsection{[Nama Divisi 2]}

\begin{itemize}[leftmargin=*, itemsep=0pt]
  \item \textbf{Kepala Divisi:} Nama kepala divisi.
  \item \textbf{Tugas:} Deskripsi tugas dan tanggung jawab divisi ini.
\end{itemize}

\section{Hak dan Wewenang}

Setiap unit kerja di \NamaPerusahaan\ memiliki hak, wewenang, serta tugas yang
jelas untuk memastikan operasional perusahaan berjalan lancar dan efisien.

% Sesuaikan sub-section dengan jabatan struktural yang ada.
\subsection{[Jabatan Utama 1]}

\begin{itemize}[leftmargin=*, itemsep=0pt]
  \item \textbf{Hak:} Hak jabatan ini.
  \item \textbf{Wewenang:} Wewenang jabatan ini.
  \item \textbf{Tugas:} Tugas utama jabatan ini.
\end{itemize}

\subsection{[Jabatan Utama 2]}

\begin{itemize}[leftmargin=*, itemsep=0pt]
  \item \textbf{Hak:} Hak jabatan ini.
  \item \textbf{Wewenang:} Wewenang jabatan ini.
  \item \textbf{Tugas:} Tugas utama jabatan ini.
\end{itemize}

\section{Lokasi \NamaPerusahaan}

\NamaPerusahaan\ berlokasi di \AlamatPerusahaan. Lokasi ini strategis dan mudah
diakses, mendukung operasional perusahaan serta memudahkan kolaborasi dengan
klien dan mitra bisnis.

\begin{itemize}[leftmargin=*, itemsep=0pt]
  \item \textbf{Nama Perusahaan:} \NamaPerusahaan
  \item \textbf{Alamat:} \AlamatPerusahaan
\end{itemize}

% Contoh penyisipan peta:
% \begin{figure}[H]
%   \centering
%   \includegraphics[width=0.85\textwidth]{assets/gambar/peta_lokasi.png}
%   \caption{Peta Lokasi \NamaPerusahaan}
%   \label{fig:peta_lokasi}
% \end{figure}

\section{Etika Profesi}

Uraikan etika dan budaya kerja yang diterapkan di \NamaPerusahaan, meliputi
standar perilaku, kedisiplinan, kerja sama tim, komunikasi, dan nilai-nilai
profesionalisme yang berlaku di lingkungan kerja instansi tersebut.

\begin{enumerate}[label=\arabic*., leftmargin=*, itemsep=0pt]
  \item \textbf{Sikap dan Kedisiplinan:} Deskripsi etika pertama.
  \item \textbf{Kerjasama Tim:} Deskripsi etika kedua.
  \item \textbf{Komunikasi yang Efektif:} Deskripsi etika ketiga.
  \item \textbf{Integritas:} Deskripsi etika keempat.
\end{enumerate}

%============================================================
% BAB 3: LANDASAN TEORI
%   Flat list per teknologi/konsep, masing-masing jadi section.
%   Sub-section untuk sub-topik teknologi jika diperlukan.
%   Sesuaikan dengan teknologi yang digunakan pada proyek Anda.
%============================================================
\chapter{LANDASAN TEORI}

Bab ini memuat teori-teori dan teknologi yang menjadi landasan dalam pelaksanaan
kerja praktik di \NamaPerusahaan. Teori yang dibahas meliputi konsep dasar,
arsitektur, kelebihan, dan implementasi umum dari teknologi-teknologi yang
digunakan dalam proyek selama kerja praktik berlangsung.

% -- Sesuaikan section berikut dengan teknologi yang digunakan ----------------

\section{[Teknologi/Konsep Pertama]}

Jelaskan konsep dasar teknologi atau framework pertama yang digunakan.
Sertakan definisi, tujuan penggunaan, dan bagaimana teknologi ini
berkontribusi pada proyek selama kerja praktik.

% Tambahkan sub-section jika teknologi ini memiliki sub-topik penting.
\subsection{[Sub-topik jika diperlukan]}

Jelaskan sub-topik pertama dari teknologi ini.

\section{[Teknologi/Konsep Kedua]}

Jelaskan konsep dasar teknologi atau framework kedua yang digunakan.
Sertakan definisi, tujuan penggunaan, jika diperlukan.

\section{[Teknologi/Konsep Ketiga]}

Jelaskan konsep dasar teknologi atau framework ketiga yang digunakan.

\subsection{[Sub-topik 3.x.1]}

Jelaskan sub-topik pertama.

\subsection{[Sub-topik 3.x.2]}

Jelaskan sub-topik kedua.

% Tambahkan section lebih lanjut sesuai jumlah teknologi yang digunakan.

%============================================================
% BAB 4: RANCANGAN DAN IMPLEMENTASI SISTEM
%   4.1 Standar Perusahaan (jika ada)
%       4.1.x sub-standar
%   4.2 Proyek [Nama Proyek 1]
%       4.2.1 Rancangan Sistem
%   4.3 Proyek [Nama Proyek 2]
%       4.3.1 Rancangan Sistem
%   (hapus section yang tidak relevan jika hanya 1 proyek)
%============================================================
\chapter{RANCANGAN DAN IMPLEMENTASI SISTEM}

Bab ini menjelaskan proses perancangan dan implementasi sistem selama kerja
praktik di \NamaPerusahaan\ dalam periode \PeriodeKP. Pembahasan mencakup
standar perusahaan dalam pengembangan perangkat lunak serta implementasi
proyek-proyek yang dikerjakan.

\section{Standar Perusahaan}

Jelaskan standar pengembangan perangkat lunak yang diterapkan di \NamaPerusahaan\
dan wajib diikuti selama pelaksanaan kerja praktik. Standar ini menjadi landasan
dalam proses pengembangan semua proyek.

% Sesuaikan sub-section dengan standar yang berlaku di perusahaan.
\subsection{[Standar 1]}

Jelaskan standar version control atau branching model yang diterapkan.

\subsubsection{[Sub-standar]}

Jelaskan detail branching model yang digunakan.

\subsubsection{[Sub-standar]}

Jelaskan standar pesan commit yang diterapkan.

\subsection{[Standar 2]}

Jelaskan standar penamaan atau konvensi database yang diterapkan.

\subsection{[Standar 3]}

Jelaskan standar struktur direktori atau arsitektur proyek yang digunakan.

\section{Proyek [Nama Proyek 1]}

Jelaskan latar belakang, tujuan, dan deskripsi aplikasi dari proyek pertama
yang dikerjakan selama kerja praktik.

\subsection{Rancangan Sistem}

Jelaskan perancangan sistem dari proyek pertama, meliputi arsitektur sistem,
desain basis data, alur data, dan komponen teknis lainnya. Sertakan diagram
yang relevan.

% Contoh penyisipan gambar arsitektur:
% \begin{figure}[H]
%   \centering
%   \includegraphics[width=0.9\textwidth]{assets/gambar/arsitektur_proyek1.png}
%   \caption{Desain Arsitektur Sistem [Nama Proyek 1]}
%   \label{fig:arsitektur_proyek1}
% \end{figure}

\subsubsection{[Sub-rancangan]}

Jelaskan detail sub-komponen rancangan pertama.

\subsubsection{[Sub-rancangan]}

Jelaskan detail sub-komponen rancangan kedua.

% Hapus section berikut jika hanya mengerjakan 1 proyek.
\section{Proyek [Nama Proyek 2]}

Jelaskan latar belakang, tujuan, dan deskripsi aplikasi dari proyek kedua
yang dikerjakan selama kerja praktik.

\subsection{Rancangan Sistem}

Jelaskan perancangan sistem dari proyek kedua, meliputi arsitektur sistem,
desain basis data, alur data, dan komponen teknis lainnya.

% Contoh penyisipan gambar arsitektur:
% \begin{figure}[H]
%   \centering
%   \includegraphics[width=0.9\textwidth]{assets/gambar/arsitektur_proyek2.png}
%   \caption{Desain Arsitektur Sistem [Nama Proyek 2]}
%   \label{fig:arsitektur_proyek2}
% \end{figure}

%============================================================
% BAB 5: PENUTUP
%   5.1 Kesimpulan
%   5.2 Saran
%       5.2.1 Saran untuk Perusahaan
%       5.2.2 Saran untuk Pihak Kampus
%       5.2.3 Saran untuk Mahasiswa Selanjutnya
%       5.2.4 Saran Pengembangan [Proyek 1] (hapus jika tidak relevan)
%       5.2.5 Saran Pengembangan [Proyek 2] (hapus jika tidak relevan)
%============================================================
\chapter{PENUTUP}

\section{Kesimpulan}

Berdasarkan pelaksanaan kerja praktik yang dilakukan di \NamaPerusahaan\ selama
periode \PeriodeKP, dapat disimpulkan sebagai berikut:

\begin{enumerate}[label=\arabic*., leftmargin=*, itemsep=0pt]
  \item Kesimpulan pertama berkaitan dengan pencapaian tujuan kerja praktik.
  \item Kesimpulan kedua berkaitan dengan hasil proyek yang dikerjakan.
  \item Kesimpulan ketiga berkaitan dengan kompetensi yang diperoleh selama
        pelaksanaan kerja praktik.
\end{enumerate}

\section{Saran}

Setelah selesai menjalani kerja praktik dan mengerjakan laporan ini, penulis
menyampaikan beberapa saran sebagai bahan masukan bagi berbagai pihak.

\subsection{Saran untuk Perusahaan (\NamaPerusahaan)}

\begin{enumerate}[label=\arabic*., leftmargin=*, itemsep=0pt]
  \item Saran bagi instansi terkait pengembangan sistem atau proses kerja
        di masa mendatang.
  \item Saran kedua untuk instansi.
\end{enumerate}

\subsection{Saran untuk Pihak Kampus (Politeknik Elektronika Negeri Surabaya)}

\begin{enumerate}[label=\arabic*., leftmargin=*, itemsep=0pt]
  \item Saran bagi program studi terkait persiapan mahasiswa sebelum kerja praktik.
  \item Saran kedua untuk program studi.
\end{enumerate}

\subsection{Saran untuk Mahasiswa Selanjutnya}

\begin{enumerate}[label=\arabic*., leftmargin=*, itemsep=0pt]
  \item Saran bagi mahasiswa yang akan melaksanakan kerja praktik di tempat
        yang sama atau bidang yang serupa.
  \item Saran kedua untuk mahasiswa selanjutnya.
\end{enumerate}

% Hapus sub-section berikut jika tidak ada saran pengembangan per proyek.
\subsection{Saran Pengembangan untuk [Nama Proyek 1]}

\begin{enumerate}[label=\arabic*., leftmargin=*, itemsep=0pt]
  \item Saran pengembangan teknis untuk proyek pertama ke depannya.
  \item Saran kedua terkait pengembangan proyek pertama.
\end{enumerate}

\subsection{Saran Pengembangan untuk [Nama Proyek 2]}

\begin{enumerate}[label=\arabic*., leftmargin=*, itemsep=0pt]
  \item Saran pengembangan teknis untuk proyek kedua ke depannya.
  \item Saran kedua terkait pengembangan proyek kedua.
\end{enumerate}


% -- Daftar Pustaka -------------------------------------------
\cleardoublepage
\addcontentsline{toc}{chapter}{DAFTAR PUSTAKA}
\renewcommand{\bibname}{DAFTAR PUSTAKA}
\bibliographystyle{apalike} % urut abjad, nama-tahun (Pasal 4.2.2 no. 7)
\bibliography{refs/references}

% -- Lampiran -------------------------------------------------
%============================================================
% LAMPIRAN
% Lampiran berisi dokumen pendukung seperti surat keterangan
% KP, logbook kegiatan, dokumentasi proyek, dan lain-lain.
%============================================================
\cleardoublepage
\appendix
\phantomsection
\addcontentsline{toc}{chapter}{LAMPIRAN}
\begin{center}
  {\FontJudulHalaman\bfseries LAMPIRAN\par}
\end{center}
\vspace{12pt}
\renewcommand{\thesection}{\Alph{section}}

Lampiran berisi dokumen-dokumen pendukung laporan kerja praktik, seperti surat
keterangan telah melaksanakan kerja praktik, logbook kegiatan harian, dokumentasi
foto atau tangkapan layar proyek, dan dokumen teknis lainnya.

% -- Lampiran Wajib (Pasal 4.2.1 ix) --------------------------
% Urutan A-F harus sama dengan pages/daftar_lampiran.tex.
\section{Rekap Monitoring Kegiatan KP yang Telah Diverifikasi}
Sisipkan rekap monitoring kegiatan KP di sini.

\section{Nilai Pembimbing Perusahaan}
Sisipkan lembar nilai dari pembimbing perusahaan di sini.

\section{Nilai Dosen Pembimbing KP}
Sisipkan lembar nilai dari dosen pembimbing KP di sini.

\section{Nilai Dosen Penguji KP}
Sisipkan lembar nilai dari dosen penguji KP di sini.

\section{Foto Mahasiswa KP di Lokasi KP}
Sisipkan foto mahasiswa di lokasi KP di sini.

\section{Biodata Penulis}
Nama, tempat dan tanggal lahir, alamat, email, serta riwayat pendidikan penulis.

% -- Lampiran Tambahan (opsional) -----------------------------
% Data pendukung, kuesioner, listing program, desain mekanik, dll.
% \section{Surat Keterangan Kerja Praktik}
% \includepdf[pages=-]{assets/lampiran/surat_kp.pdf}

% \section{Logbook Kegiatan Harian}
% Isi logbook di sini.

% \section{Dokumentasi Proyek}
% Isi dokumentasi di sini.


%============================================================
\end{document}
%============================================================
